\chapter{ฟิสิกส์ทุเรียนและพลศาสตร์การเจริญเติบโต}
\label{chap:ch02}

\section*{แผนบริหารการสอนประจำบทที่ 2}
\addcontentsline{toc}{section}{แผนบริหารการสอนประจำบทที่ 2}

\noindent\textbf{หัวข้อเนื้อหาประจำบท}
\begin{enumerate}
    \item กฎฟิสิกส์และเทอร์โมไดนามิกส์ในระบบต้นทุเรียน
    \item ชีวกลศาสตร์ของขั้วผลและการกระจายแรงโน้มถ่วง
    \item จลนพลศาสตร์การสะสมเนื้อแห้งและเกณฑ์การเก็บเกี่ยว
\end{enumerate}

\noindent\textbf{วัตถุประสงค์เชิงพฤติกรรม}
\begin{enumerate}
    \item อธิบายหลักการ ทฤษฎี และสถาปัตยกรรมเทคโนโลยีประจำบทที่ 2 ได้อย่างถูกต้อง
    \item ประยุกต์ใช้เครื่องมือ อุปกรณ์เซนเซอร์ และระบบปัญญาประดิษฐ์เพื่อการจัดการสวนทุเรียนได้อย่างมีประสิทธิภาพ
    \item วิเคราะห์ สรุปผลข้อมูล และแก้ไขปัญหาเชิงฟิสิกส์เกษตรและคุณภาพดินได้อย่างเป็นระบบ
\end{enumerate}

\noindent\textbf{การวัดและประเมินผล}
\begin{table}[H]
\centering\small
\begin{tabularx}{\textwidth}{p{0.55\textwidth} p{0.25\textwidth} c}
\toprule
\textbf{เกณฑ์การประเมินผล} & \textbf{เครื่องมือวัดผล} & \textbf{สัดส่วนคะแนน} \\
\midrule
1. ทักษะการปฏิบัติงานและการใช้อุปกรณ์ & แบบประเมินทักษะภาคสนาม & 40\% \\
2. รายงานข้อมูลและการวิเคราะห์ผล & การตรวจสมุดบันทึกและดาต้าชีท & 30\% \\
3. แบบฝึกหัดและคำถามท้ายบท & แบบทดสอบท้ายบทเรียน & 20\% \\
4. การมีส่วนร่วมและความรับผิดชอบ & การเข้าเรียนและการตรงต่อเวลา & 10\% \\
\midrule
\textbf{รวมทั้งสิ้น} & & \textbf{100\%} \\
\bottomrule
\end{tabularx}
\end{table}
\newpage

% ==========================================================
% เนื้อหาบทเรียนประจำบทที่ 2
% ==========================================================

\section{กฎฟิสิกส์และเทอร์โมไดนามิกส์ในระบบต้นทุเรียน}
\index{กฎฟิสิกส์และเทอร์โมไดนามิกส์ในระบบต้นทุเรียน}

การศึกษาและทำความเข้าใจเกี่ยวกับกฎฟิสิกส์และเทอร์โมไดนามิกส์ในระบบต้นทุเรียน (Thermodynamics and Heat Budget in Durian Canopy) มีความสำคัญอย่างยิ่งในระบบเกษตรอัจฉริยะ โดยมีรายละเอียดสาระสำคัญดังนี้

ต้นทุเรียนจัดเป็นระบบอุณหพลศาสตร์เปิด (Open thermodynamic system) ที่มีการแลกเปลี่ยนทั้งมวลสาร (น้ำ, คาร์บอนไดออกไซด์, สารอาหาร) และพลังงานกับสิ่งแวดล้อมรอบตัว

สมดุลความร้อนของทรงพุ่มทุเรียน (Canopy Energy Budget) อธิบายได้ด้วยสมการ:
\[ R\_n = H + \lambda E + G + S \]
โดยที่ $R_n$ คือรังสีสุทธิจากดวงอาทิตย์ (Net radiation), $H$ คือฟลักซ์ความร้อนสัมผัส (Sensible heat flux), $\lambda E$ คือฟลักซ์ความร้อนแฝงจากการคายระเหยน้ำ (Latent heat flux of transpiration), $G$ คือฟลักซ์ความร้อนลงสู่ผิวดิน (Soil heat flux), และ $S$ คืออัตราการสะสมความร้อนในชีวมวลของลำต้นและผล

ในวันที่อากาศร้อนจัด หากต้นทุเรียนขาดน้ำ กระบวนการคายน้ำเพื่อระบายความร้อนแฝง (Transpirational cooling) จะหยุดชะงัก ส่งผลให้อุณหภูมิใบพุ่งสูงเกิน 40$^{\circ}$C และเกิดความเสียหายต่อระบบสังเคราะห์แสง (Photoinhibition)

\begin{conceptbox}[มโนทัศน์สำคัญทางฟิสิกส์เกษตร]
กฎการอนุรักษ์พลังงานและสมดุลความร้อนของทรงพุ่มเป็นตัวแปรสำคัญที่กำหนดอัตราการคายน้ำและการสังเคราะห์แสงของต้นทุเรียน การควบคุมปัจจัยแวดล้อมจึงต้องสอดคล้องกับหลักการทางฟิสิกส์
\end{conceptbox}

\section{ชีวกลศาสตร์ของขั้วผลและการกระจายแรงโน้มถ่วง}
\index{ชีวกลศาสตร์ของขั้วผลและการกระจายแรงโน้มถ่วง}

การศึกษาและทำความเข้าใจเกี่ยวกับชีวกลศาสตร์ของขั้วผลและการกระจายแรงโน้มถ่วง (Biomechanics of Durian Fruit Stalk and Mass Distribution) มีความสำคัญอย่างยิ่งในระบบเกษตรอัจฉริยะ โดยมีรายละเอียดสาระสำคัญดังนี้

ผลทุเรียนเมื่อพัฒนาจนเติบโตเต็มที่อาจมีน้ำหนักตั้งแต่ 2.5 ถึง 6.0 กิโลกรัม การแขวนลอยของผลบนกิ่งต้องอาศัยโครงสร้างทางชีวกลศาสตร์ (Biomechanics) ที่แข็งแรงเป็นพิเศษ

ขั้วผลทุเรียน (Pedicel) ประกอบด้วยเนื้อเยื่อไซเลมและเส้นใยสเกลอเรนไคมา (Sclerenchyma fibers) ที่มีความทนทานต่อแรงดึง (Tensile stress) สูงมาก แรงกระทำต่อขั้วผลคำนวณตามกฎของนิวตัน:
\[ F = m(g + a) \]
โดยที่ $m$ คือมวลของผลทุเรียน, $g$ คือความเร่งเนื่องจากแรงโน้มถ่วง ($9.81$ m/s$^2$), และ $a$ คือความเร่งพลศาสตร์ที่เกิดจากแรงลมพัดโยกกิ่ง หากแรงลมกระโชก ($a$) มีค่าสูง แรงดึงรวมอาจเกินขีดจำกัดความเหนียวของเนื้อเยื่อขั้วผล ทำให้ผลร่วงหล่นก่อนกำหนด การเข้าใจหลักฟิสิกส์นี้ช่วยในการออกแบบการโยงกิ่งและผูกเชือกพยุงลูกทุเรียนอย่างถูกต้องตามหลักวิศวกรรม

\section{จลนพลศาสตร์การสะสมเนื้อแห้งและเกณฑ์การเก็บเกี่ยว}
\index{จลนพลศาสตร์การสะสมเนื้อแห้งและเกณฑ์การเก็บเกี่ยว}

การศึกษาและทำความเข้าใจเกี่ยวกับจลนพลศาสตร์การสะสมเนื้อแห้งและเกณฑ์การเก็บเกี่ยว (Dry Matter Kinetics and Harvesting Standards) มีความสำคัญอย่างยิ่งในระบบเกษตรอัจฉริยะ โดยมีรายละเอียดสาระสำคัญดังนี้

คุณภาพความอร่อยของทุเรียนขึ้นอยู่กับเปอร์เซ็นต์เนื้อแห้ง (Dry Matter Percentage; \%DM) ซึ่งสะท้อนการสะสมแป้งและน้ำตาลในเนื้อทุเรียน

เกณฑ์มาตรฐานสินค้าเกษตร (มกษ. 9003-2556) กำหนดเปอร์เซ็นต์น้ำหนักเนื้อแห้งขั้นต่ำสำหรับการเก็บเกี่ยว:
- พันธุ์กระดุมทอง: ไม่น้อยกว่า 27\%
- พันธุ์ชะนี: ไม่น้อยกว่า 30\%
- พันธุ์พวงมณี: ไม่น้อยกว่า 30\%
- พันธุ์หมอนทอง: ไม่น้อยกว่า 32\%

การสะสมเนื้อแห้งสามารถจำลองได้ด้วยสมการซิกมอยด์ (Sigmoidal curve) ซึ่งอัตราการสะสมแป้งจะเพิ่มขึ้นสูงสุดในช่วง 90 ถึง 110 วันหลังดอกบานเต็มที่ (Days After Anthesis; DAA)

\section{ขั้นตอนและวิธีปฏิบัติการทดลอง}
\index{ขั้นตอนการทดลองประจำบทที่ 2}

\subsection{การวัดรังสีแสงอาทิตย์และอุณหภูมิผิวผลทุเรียนด้วยกล้องถ่ายภาพความร้อน}
\begin{sensorbox}[การวัดรังสีแสงอาทิตย์และอุณหภูมิผิวผลทุเรียนด้วยกล้องถ่ายภาพความร้อน]
1. ติดตั้งเซนเซอร์วัดรังสีดวงอาทิตย์ (Pyranometer) เหนือทรงพุ่มทุเรียน บันทึกค่า Irradiance ($W/m^2$)
2. ใช้กล้องถ่ายภาพความร้อนอินฟราเรด (Thermal Imaging Camera) บันทึกภาพผลทุเรียนในทิศตะวันออกและทิศตะวันตก
3. ตรวจวัดอุณหภูมิผิวเปลือกทุเรียน (Surface temperature) เปรียบเทียบกับอุณหภูมิอากาศรอบข้าง
4. วิเคราะห์จุดเสี่ยงการเกิดอาการแดดเผา (Sunburn) บนเปลือกทุเรียนเมื่ออุณหภูมิผิวสูงเกิน 42$^{\circ}$C
\end{sensorbox}

\subsection{การตรวจสอบเปอร์เซ็นต์น้ำหนักเนื้อแห้งของทุเรียนด้วยวิธีอบความร้อน}
\begin{sensorbox}[การตรวจสอบเปอร์เซ็นต์น้ำหนักเนื้อแห้งของทุเรียนด้วยวิธีอบความร้อน]
1. สุ่มกรีดเจาะเนื้อทุเรียนบริเวณกึ่งกลางพู น้ำหนักสดประมาณ 10-15 กรัม
2. หั่นเนื้อทุเรียนเป็นแผ่นบางความหนาประมาณ 1-2 มิลลิเมตร ชั่งน้ำหนักสดเริ่มต้น ($W_1$) ด้วยเครื่องชั่งละเอียด 4 ตำแหน่ง
3. อบเนื้อทุเรียนในตู้อบลมร้อนที่อุณหภูมิ 105$^{\circ}$C เป็นเวลา 4 ชั่วโมง หรือจนกระทั่งน้ำหนักคงที่
4. นำออกมาผึ่งในโถดูดความชื้น (Desiccator) จนเย็น ชั่งน้ำหนักแห้ง ($W_2$)
5. คำนวณเปอร์เซ็นต์เนื้อแห้งตามสมการ: $\%DM = (W_2 / W_1) \times 100$
\end{sensorbox}

\subsection{การทดสอบแรงดึงของขั้วผลทุเรียนด้วยเครื่องวัดแรงดิจิทัล}
\begin{sensorbox}[การทดสอบแรงดึงของขั้วผลทุเรียนด้วยเครื่องวัดแรงดิจิทัล]
1. ยึดขั้วผลทุเรียนเข้ากับแท่นจับทดสอบแรงดึงดิจิทัล (Digital Force Gauge)
2. เพิ่มแรงดึงในแนวดิ่งอย่างสม่ำเสมอด้วยความเร็วคงที่ 5 มม./นาที
3. บันทึกค่าแรงดึงสูงสุด (Ultimate Tensile Force; $N$) ณ จุดที่ขั้วผลเริ่มปริขาด
4. เปรียบเทียบค่าความเหนียวของขั้วผลในทุเรียนที่ได้รับแคลเซียม-โบรอนตามโปรแกรมปุ๋ยกับชุดควบคุม
\end{sensorbox}

\section{ตารางบันทึกผลการทดลองและการอภิปรายผล}
\begin{table}[H]
\centering\small
\caption{ตารางบันทึกผลการตรวจวัดและข้อมูลพารามิเตอร์ประจำบทที่ 2}
\label{tab:ch02_datasheet}
\begin{tabularx}{\textwidth}{p{0.3\textwidth} p{0.3\textwidth} X}
\toprule
\textbf{พารามิเตอร์ / การทดสอบ} & \textbf{ค่าที่ตรวจวัดได้} & \textbf{การแปลผลและการวิเคราะห์} \\
\midrule
จุดตรวจวัดที่ 1 (บริเวณดินดอน) & ค่าความชื้นอยู่ในเกณฑ์ปกติ & ปริมาณน้ำเพียงพอต่อการเจริญ \\
จุดตรวจวัดที่ 2 (บริเวณที่ลุ่ม) & มีการสะสมความชื้นสูง & ควรชะลอการให้น้ำเพื่อป้องกันรากเน่า \\
สถานะสัญญาณโครงข่าย & RSSI: -85 dBm, SNR: +8 dB & คุณภาพสัญญาณ LoRaWAN ยอดเยี่ยม \\
\bottomrule
\end{tabularx}
\end{table}

\section{คำถามทบทวนและแบบฝึกหัดท้ายบท}
\begin{enumerate}
    \item จงอธิบายผลกระทบของอุณหภูมิสูงและรังสีแสงอาทิตย์จัดต่อการทำงานของปากใบ (Stomatal Conductance) ในต้นทุเรียน
    \item เหตุใดการตรวจวัดเปอร์เซ็นต์เนื้อแห้ง (\%DM) จึงเป็นตัวชี้วัดความสุกแก่ของทุเรียนที่แม่นยำกว่าการนับอายุวันเพียงอย่างเดียว
    \item หากเกิดพายุฤดูร้อนที่มีความเร็วลมกระโชก 50 km/h จงคำนวณแรงพลศาสตร์ที่กระทำต่อผลทุเรียนน้ำหนัก 3.5 kg
\end{enumerate}
